\chapter{Conteúdo do Projeto}

\section{Bases de Dados (Datasets)}

As instâncias para teste e validação dos algoritmos devem ser extraídas das seguintes fontes acadêmicas:

\begin{itemize}
\item TSPLIB95: \url{http://www.iwr.uni-heidelberg.de/groups/comopt/software/TSPLIB95/}
\item World TSP (University of Waterloo): 
\url{https://www.math.uwaterloo.ca/tsp/world/countries.html}
\item Repositório do professor Vinicius de Novaes com alguns exemplos organizados das fontes anteriores: \url{https://github.com/viniciusdenovaes/TSPLib-APS_IA}
\end{itemize}


\section{Desenvolvimento e Conceitos Computacionais}

O trabalho será dividido em três etapas principais de implementação:



\subsection{Busca em Espaço de Estados (Solução Ótima)}

Nesta etapa, o grupo deve implementar um algoritmo que faça uma busca exaustiva no Espaço de Estados (podendo ser uma busca em Profundidade ou Largura).

\begin{description}
    \item[Conceito:] Busca em Espaço de Estados (Solução Ótima) Nesta etapa, o grupo deve implementar um algoritmo que faça uma busca exaustiva no Espaço de Estados (podendo ser uma busca em Profundidade ou Largura).
    \item[Análise:] Deve-se discutir a explosão combinatória de complexidade O(n!).Neste caso teremos n! possibilidades, sendo n o número de vértices no grafo, mostrando como o tempo de execução cresce exponencialmente com o aumento do número de cidades (n).

\end{description}








\subsection{Branch and Bound (Solução Ótima Otimizada)}

A aplicação da técnica de Branch and Bound visa encontrar \textbf{a mesma} solução ótima do passo anterior, porém de forma mais eficiente.

\begin{description}
    \item[Conceito:] A técnica utiliza a árvore de busca onde cada nó representa uma solução parcial. O "Branch" (ramificação) gera os filhos de cada vértice, enquanto o "Bound" (limite) calcula um limite inferior de custo para cada ramo.

    Poda (Pruning): Se o limite inferior de um ramo for maior que a melhor solução já encontrada (upper bound), o ramo é descartado ("podado"), evitando a exploração de caminhos inúteis e reduzindo o tempo de processamento em relação à busca exaustiva simples.

    \item[Complexidade:] Deve-se discutir a melhora no tempo de execução em relação ao anterior.

\end{description}






\subsection{Heurística (Solução Aproximada)}

O grupo deve implementar uma técnica que não garante a otimalidade, mas entrega uma solução "boa o suficiente" em tempo polinomial. Para este item, sugere-se o uso de Soluções iniciais Greedy Algoritmos Genéticos, Colônia de Formiga ou Busca Local (ex: 2-opt). O grupo pode escolher qualquer um.

\begin{description}
    \item[Conceito:] Diferente dos métodos exatos, as heurísticas focam na eficiência computacional. 
    Uma solução Greedy, por exemplo, pode escolher apenas as menores rotas disponíveis até que uma solução seja formada.
    No caso de Algoritmos Genéticos, utiliza-se conceitos de evolução biológica (população, cruzamento, mutação e seleção natural) para iterativamente melhorar as rotas. 

    \item[Análise:] O foco aqui é demonstrar que, para instâncias grandes (onde os passos 1 e 2 falham por tempo/memória), a heurística consegue fornecer resultados viáveis rapidamente.

\end{description}





